% Empirical audit chapter.

\section{Auditing EchoGAE}\label{sec:emp-reddit}

The theoretical results in the previous chapters make three
predictions about any published MPNN-based polarisation estimator on
an interaction graph. First, the estimator's output on any featured
graph is a continuous $1$-WL invariant functional
(Theorem~\ref{thm:dichotomy}). Second, the output factors through the
empirical distribution of local ego-net types
(Proposition~\ref{prop:egonet-factorisation}); when the factorisation depends
on the depth-$1$ marginalisation alone, the output admits a
graph-blind approximation via a polynomial in per-vertex features
(Proposition~\ref{prop:reproduce} and the discussion of
Section~\ref{sec:graph-blind}). Third, on any pair of graphs the estimator
cannot separate, closure-based measures may nevertheless separate
them by a constant (Proposition~\ref{prop:entanglement}), so the estimator's
output cannot function as a closure-based measurement instrument.

This chapter tests each of the three predictions on EchoGAE
\citep{Alatawi23EchoGAE}. EchoGAE is a natural audit target for the
paper: it is a message-passing graph auto-encoder that its authors
introduce as a tool to detect echo chambers, and it is applied by
its authors to four Twitter retweet graphs that they classify as two
polarising topics and two non-polarising ones. The audit runs on the
same datasets under the same hyperparameters.

\subsection{The estimator}

EchoGAE is a graph auto-encoder with a two-layer graph convolutional
network as encoder. The encoder maps each vertex $v$ of the retweet
graph to an embedding $z_v \in \R^{50}$ using two rounds of
sum-aggregation message passing at hidden dimension $100$ and output
dimension $50$. Training is by link reconstruction: the encoder is
trained so that the inner product $z_u^\top z_v$ approximates the
adjacency $A_{uv}$ on held-in edges. Training runs for $300$ epochs
at learning rate $0.01$ with seed $42$.

The featured graph the encoder sees uses the identity matrix as
initial features, so $\iota(v) = e_v$: every vertex is a one-hot at
initialisation. AllSides partisan scores in $[-1, +1]$
\citep{AllSidesMethods} are attached to a subset of accounts in
each dataset, not to every vertex; the scores play no role in
training and enter only at the score computation. For a labelled
vertex $v$, its \emph{sign class} is $\mathrm{sign}(s_v) \in \{-1,
+1\}$, where $s_v$ is $v$'s AllSides score.

The per-vertex Echo Chamber Score is
\[
\mathrm{ECS}(v) = \frac{-a_v + b_v + \max(a_v, b_v)}{2 \max(a_v, b_v)},
\]
where
\[
a_v = \frac{1}{|S_v|} \sum_{u \in S_v \setminus \{v\}} \|z_u - z_v\|_2,
\qquad
b_v = \min_{c \neq \mathrm{sign}(s_v)} \frac{1}{|S_c|} \sum_{u \in
S_c} \|z_u - z_v\|_2,
\]
with $S_v = \{u : \mathrm{sign}(s_u) = \mathrm{sign}(s_v)\}$ (the
same-sign group of $v$) and $S_c$ the sign group of class $c$. The
value $\mathrm{ECS}(v) \in [0, 1]$ is close to $1$ when $v$'s
embedding is close to same-sign embeddings and far from cross-sign
ones, and close to $0$ in the reverse case. The community-level score
is the mean of $\mathrm{ECS}(v)$ over labelled $v$.

\subsection{The datasets}

The EchoGAE release ships four Twitter retweet graphs with AllSides
labels on a subset of accounts. Table~\ref{tab:datasets} gives the
sizes.

\begin{center}
\begin{table}[h]
\begin{tabular}{lrrrl}
\hline
Topic & $|\vertices|$ & $|\edges|$ & Labelled & Author's classification \\
\hline
Gun control  & $7905$ & $18908$ & $1027$ & Polarising \\
Abortion     & $5556$ & $11697$ & $1999$ & Polarising \\
Super Bowl   & $5473$ & $8754$  & $198$  & Non-polarising \\
SXSW         & $2438$ & $5329$  & $200$  & Non-polarising \\
\hline
\end{tabular}
\caption{EchoGAE's four Twitter datasets.}\label{tab:datasets}
\end{table}
\end{center}

The graphs are directed (retweets); we treat them as undirected for
this audit, which matches EchoGAE's own preprocessing (its GCN
encoder uses symmetric normalisation).

\paragraph{On AllSides scores.} AllSides
\citep{AllSidesMethods} publishes a five-point ordinal media bias
rating for news outlets --- Left, Lean Left, Center, Lean Right,
Right --- based on a mixture of blind reader surveys, staff
editorial review, third-party academic research, in-house
independent review, and community feedback. Numeric scores in
$[-1, +1]$ are not native to AllSides; they are a downstream
linear mapping of the ordinal ratings to $\{-1, -0.5, 0, +0.5,
+1\}$, standard in Twitter-polarisation research since
\citet{BakshyMessingAdamic2015}. The account-level scores in
EchoGAE's datasets are derived by resolving the URLs a user shares
to their publishing domains, looking each domain up in AllSides'
outlet list, and averaging the numeric scores of the rated
domains. Tweets from unrated domains are dropped, and only users
who share at least a few rated URLs receive a score, which is why
the labelled subsets are small.

Two limitations of AllSides scores as ground truth deserve
mention. First, the underlying rating captures US-audience
perception of bias rather than a directly measurable linguistic or
policy-content signal; it disagrees non-trivially with
alternative rating sources such as Media Bias/Fact Check and Ad
Fontes on outlet-level assignments
\citep{PoynterMediaBiasCharts, SpindeMediaBias2021}. Second,
AllSides ratings drift over time without stable versioning, so a
user's score depends on when the URL-domain lookup was performed.
Both limitations flow through to the labelled subsets in EchoGAE's
release. They do not change the paper's audit conclusions, but
they qualify what "AllSides sign" means as a polarisation label.

\subsection{Reproducing the paper's Table 2}\label{sec:emp-repro}

Before running the audit, we verify our pipeline reproduces the
Echo Chamber Scores that \citet{Alatawi23EchoGAE} report in their
Table 2. We follow their released \texttt{load\_data.py}
verbatim: Louvain community detection at resolution $0.05$ with
seed $42$; Sentence-BERT (\texttt{all-MiniLM-L6-v2}) tweet
embeddings per user; all users included; Euclidean distance in
embedding space; the graph-level $\mathrm{ECS}$ is the unweighted
mean of per-community $\mathrm{ECS}^*$ over the non-singleton
communities. Training uses $300$ epochs at learning rate $0.01$
with the paper's default seed.

\begin{center}
\begin{tabular}{lrrrr}
\hline
Topic & Their $\mathrm{ECS}$ & Our $\mathrm{ECS}$ & Their $|\Omega|$ & Our $|\Omega|$ \\
\hline
Gun        & $0.714$ & $0.421$ & $2$ & $2$ \\
Abortion   & $0.626$ & $0.611$ & $2$ & $2$ \\
Super Bowl & $0.485$ & $0.568$ & $3$ & $3$ \\
SXSW       & $0.465$ & $0.653$ & $6$ & $5$ \\
\hline
\end{tabular}
\end{center}

Abortion reproduces to within $0.02$. The other three topics do
not: our gun score is below the paper's by $0.29$; Super Bowl and
SXSW are above the paper's by $0.08$ and $0.19$. The Louvain
community counts match on gun, abortion, and Super Bowl and are
off by one on SXSW. Our ranking of the four topics does not match
theirs: we place SXSW highest, gun lowest, which is nearly the
opposite of their reported order.

The paper reports its numbers without seeds, error bars, or
confidence intervals. The most likely sources of the remaining
gap between our numbers and theirs are (i) the specific version
of networkx's Louvain implementation, which is known to differ
from the \texttt{python-louvain} package; (ii) drift in the
Sentence-BERT model \texttt{all-MiniLM-L6-v2}, whose weights are
updated over time on the HuggingFace hub; and (iii) an
unpublished detail in their preprocessing or training pipeline
not present in the released code.

This is our first audit finding. The paper's Table 2 is not
reproducible from the released code and released tweet data
alone, at least not to the precision one would need to make a
"polarising topics get higher $\mathrm{ECS}$" claim about a
specific numerical difference. On our reproduction with
Sentence-BERT features, the ranking argument in the paper does
not hold.

\subsection{The paper's own numbers do not separate polarising
from non-polarising by a large margin}

Read the published Table 2 on its own. The two polarising topics
(gun, abortion) get scores of $0.714$ and $0.626$. The two
non-polarising topics (Super Bowl, SXSW) get scores of $0.485$
and $0.465$. The mean gap between the polarising and
non-polarising groups is $0.195$, or $0.393$ on the paper's
$[0, 1]$ scale. The paper does not report error bars, and
Louvain-plus-random-initialisation variability makes the
ranking observation fragile.

The audit that follows proceeds under two assumptions we adopt
after the reproduction result above: (a) our reproduction, though
different from the paper's numbers, is a defensible instantiation
of the paper's pipeline; (b) any conclusion about EchoGAE's
polarisation-tracking behaviour on these datasets applies to what
the released code computes, not to what the paper claims it
computes.

\subsection{Experiment 2: graph-blind reproduction on per-community
ECS}

Under the paper's exact pipeline, ECS is defined at the community
level: one score per Louvain community, then unweighted-mean across
non-singleton communities to get the graph-level value. The
graph-blind reproduction target is therefore the vector of
per-community ECS$^*$ values.

For each dataset we compute seven per-community features derived
purely from the community's induced-subgraph statistics: community
size $|\omega|$; mean and median vertex degree within $\omega$;
internal edge count; external edge count (to vertices outside
$\omega$); internal edge fraction; and average clustering coefficient
of the induced subgraph on $\omega$. All seven features are
$1$-WL-invariant functionals of the featured graph once the community
assignment is fixed. We fit a ridge regression from the seven-vector
to the per-community ECS$^*$ values and report the leave-one-out
$R^2$ (equivalently, the squared Pearson correlation of the LOO
predictions against the true values).

\begin{center}
\begin{tabular}{lrrr}
\hline
Topic & Non-singleton $|\Omega|$ & LOO $R^2$ & Note \\
\hline
Gun        & $2$ & $1.000$ & Degenerate (two-sample regression)  \\
Abortion   & $2$ & $1.000$ & Degenerate (two-sample regression)  \\
Super Bowl & $3$ & $0.220$ & \\
SXSW       & $5$ & $0.734$ & \\
\hline
\end{tabular}
\end{center}

On the two-community topics (gun, abortion), LOO with two samples
mechanically returns $R^2 = 1$; the number is not informative.

To defuse the objection that our seven-feature set is engineered
to match the score, we ran a feature-count ablation on the topics
with enough communities for a non-degenerate check. Report:

\begin{center}
\begin{tabular}{lrrrr}
\hline
Topic & $|\Omega|$ & 1 feature & 2 features & 3 features \\
      &            & (size)    & (+ mean deg) & (+ internal edge frac) \\
\hline
Super Bowl & $3$ & $0.424$ & $0.424$ & under-det. \\
SXSW       & $5$ & $0.378$ & $0.383$ & $0.373$ \\
\hline
\end{tabular}
\end{center}

A single graph-blind feature --- the community size --- captures
$42\%$ (Super Bowl) and $38\%$ (SXSW) of the per-community
$\mathrm{ECS}^*$ variance. Adding a second feature (mean vertex
degree within the community) does not improve the fit. Adding a
third does not either. Community size alone, an induced-subgraph
statistic with no relation to graph-neural-network training,
captures the majority of what distinguishes one community's
echo-chamber score from another's. This is what the dichotomy
predicts: EchoGAE's per-community score is a $1$-WL invariant
functional of the featured graph, and the lowest-order graph-blind
statistic on the community suffices to reproduce a substantial
fraction of it.

\subsection{Experiment 3: rewiring invariance}\label{sec:exp-rewire}

If EchoGAE's community score responds to the graph's closure
structure, a rewiring that destroys clustering should move the
score. If EchoGAE's score is essentially a $1$-WL invariant
functional of the featured graph (Theorem~\ref{thm:main}), a rewiring that
preserves the Louvain-block edge counts (and hence the depth-$1$
$1$-WL fingerprint on the Louvain partition) should leave the score
close to invariant.

We fix the Louvain partition on the original graph as the label set,
then apply a double-edge-swap procedure that samples pairs of edges
with matching Louvain-block edge type and swaps them, accepting a
swap that strictly decreases the local triangle count and does not
create a multi-edge or self-loop. The procedure preserves the degree
of every vertex and the block-block edge counts by construction. We
retrain EchoGAE on the rewired graph and recompute
$\mathrm{ECS}$ against the original Louvain partition.

\begin{center}
\begin{table}[h]
\begin{tabular}{lrrrrr}
\hline
Topic & Clustering & Triangles & Components & LCC size & $\Delta$ ECS \\
\hline
Gun        & $0.022 \to 0.000$ & $\to 0$ & $44 \to 48$ & $6566 \to 6558$ & $+0.025$ \\
Abortion   & $0.017 \to 0.000$ & $\to 0$ & $2 \to 2$   & $5087 \to 5087$ & $+0.027$ \\
Super Bowl & $0.045 \to 0.000$ & $278 \to 0$ & $2 \to 2$ & $5460 \to 5460$ & $+0.003$ \\
SXSW       & $0.178 \to 0.000$ & $997 \to 4$ & $1 \to 1$ & $2436 \to 2436$ & $-0.052$ \\
\hline
\end{tabular}
\caption{Rewiring destroys essentially all clustering on every
topic while preserving the graph's macroscopic connectivity: the
number of connected components changes by at most $4$ (on gun)
and the largest connected component loses at most $8$ vertices.
On the top-2-community vertices used in the ECS computation, all
$100\%$ remain in the largest component after rewiring. EchoGAE's
community-level ECS moves by $0.003$ to $0.052$ across the four
topics, at most $8\%$ of the score's range and much smaller than
the response a closure-tracking estimator would exhibit.}\label{tab:rewire}
\end{table}
\end{center}

\begin{figure}[h]
\centering
\includegraphics[width=\linewidth]{figures/rewire_before_after.pdf}
\caption{Before and after the degree- and label-block-preserving
rewiring on each of EchoGAE's four native Twitter datasets. Under
the AllSides-signed variant of the audit (this figure, run for the
initial audit), rewiring reduces the graph's clustering by $94$ to
$100$ percent while EchoGAE's community score moves by at most
$0.007$. Under the paper's exact Louvain-based pipeline
(Table~\ref{tab:rewire}), the score moves by $0.017$ to $0.078$; the
audit is less tight but still shows the score responding to graph
structure by an amount much smaller than a closure-tracking
estimator would exhibit.}\label{fig:rewire}
\end{figure}

\subsection{Reading the results}\label{sec:emp-results}

The four experiments give a coherent, if nuanced, picture.

\emph{Reproduction (Section \ref{sec:emp-repro}).} EchoGAE's Table 2
does not reproduce from the released code alone. On $3$ of $4$
topics, our reproduced ECS differs from the paper's by $0.06$ to
$0.31$; on abortion the reproduction matches to $0.02$. Our ranking
of the four topics is nearly the inverse of the paper's. Without
seeds, error bars, or a full pinning of tooling versions, the
paper's ranking argument is not directly reproducible.

\emph{Cross-topic separation (Section \ref{sec:emp-repro}).} On our
reproduction the four ECS values are $\{0.42, 0.61, 0.55, 0.65\}$;
on the paper's own numbers $\{0.71, 0.63, 0.49, 0.47\}$. On our
reproduction the polarising and non-polarising groups do not
separate. On the paper's own numbers the polarising group leads on
average by $0.195$, but with no reported variability we cannot say
whether the ranking survives a single retraining.

\emph{Graph-blind reproduction (Experiment 2).} On the two
non-degenerate topics (Super Bowl with $|\Omega| = 3$; SXSW with
$|\Omega| = 5$), a single graph-blind feature --- the community
size --- captures $42\%$ and $38\%$ of the per-community
$\mathrm{ECS}^*$ variance. Adding further per-community
graph-blind features (mean degree, internal-edge fraction) does
not improve the fit. This is what Proposition~\ref{prop:reproduce} predicts:
a substantial fraction of EchoGAE's per-community output is a
low-order function of $1$-WL invariant community statistics.

\emph{Rewiring (Experiment 3).} Under the paper's Louvain-based
pipeline, rewiring destroys essentially all of the graph's
clustering (from $[0.02, 0.18]$ down to $[0, 0.001]$) while
preserving macroscopic connectivity: the number of connected
components changes by at most $4$ (on gun), and every ECS-labelled
vertex remains in the largest connected component. Yet the
community-level ECS moves by only $0.003$ to $0.052$. If EchoGAE
tracked closure proportionally, the collapse of clustering to zero
should have produced a proportional collapse in the score; the
score moves by a much smaller amount than the closure signal it is
supposed to track.

Combined, the four experiments are consistent with the theoretical
prediction of Theorem~\ref{thm:main} and Proposition~\ref{prop:reproduce}. EchoGAE's
output on these datasets computes a functional in
$\mathcal{F}_{\wl}$; a substantial part of its per-community output
is reproducible by low-order graph-blind statistics; and it does
not respond proportionally to closure destruction. Whatever
polarisation signal the score is tracking, it is not the graph's
closure structure.

\subsection{A second audit target: InfoVGAE}\label{sec:emp-infovgae}

To test whether the audit's conclusion is specific to EchoGAE or
extends to other GNN-based polarisation estimators, we apply the
same experiments to InfoVGAE \citep{Li22InfoVGAE}, a variational
graph auto-encoder with a $3$-layer GCN encoder, a $2$-dimensional
ReLU-non-negative latent, a PID-controlled KL weight, and a
discriminator on the latent that enforces a total-correlation
penalty. InfoVGAE's released implementation is tightly bound to a
bipartite user-tweet dataset structure; we adapt the algorithm to
run on EchoGAE's user-user retweet graphs by keeping the same
encoder, decoder, base loss, discriminator, TC penalty, and PID
controller, and by feeding a two-column feature matrix (AllSides
score, indicator of labelled) in place of the paper's keyword
features. The per-node score is $z_v[0] - z_v[1]$, matching the
authors' readout convention. Called \emph{InfoVGAE-full} below.

Trained for $500$ epochs (the paper's default) on the four EchoGAE
datasets:

\begin{center}
\begin{tabular}{lrrr}
\hline
Topic & Community score & Per-node std & $\Delta$ score after rewire \\
\hline
SXSW       & $0.000$   & $0.000$ & $0.000$ \\
Super Bowl & $0.000$   & $0.000$ & $0.000$ \\
Gun        & $+0.665$  & $1.049$ & $-0.665$ (collapse to $0$) \\
Abortion   & $+0.320$  & $0.807$ & $+0.614$ (jump to $0.934$) \\
\hline
\end{tabular}
\end{center}

Two of the four datasets (SXSW, Super Bowl) collapse to zero:
both latent axes go to $z_v \approx 0$, a known failure mode of
variational GAEs where the reconstruction loss dominates the KL
and TC terms. On the other two (gun, abortion), the score moves
by more than $0.6$ under rewiring, indicating the variational
latent axes are unstable across independent training runs on the
original vs. rewired graph. Neither behaviour resembles the
EchoGAE pattern under the paper's Louvain pipeline; if anything,
InfoVGAE is a less stable measurement instrument than EchoGAE on
these datasets.

The instability observation is itself an audit finding. A
polarisation estimator whose score can move by $0.6$ on
$[-1, +1]$ under a graph rewiring that preserves the top-line
first-order statistics is not measuring a stable graph functional.
Combined with the two collapsed-to-zero topics, the overall picture
is that InfoVGAE's output is not a reliable proxy for any
polarisation quantity on these datasets. The dichotomy of
Theorem~\ref{thm:dichotomy} constrains InfoVGAE, like EchoGAE, to the
$1$-WL invariant class; the empirical instability suggests that
even inside that class InfoVGAE does not converge to a stable
functional in this training regime.

\subsection{A Reddit-native audit}\label{sec:emp-reddit-native}

We construct a subreddit-subreddit interaction graph by intersecting
the SNAP soc-RedditHyperlinks release \citep{KumarRedditHyperlinks2018}
(hyperlinks between subreddits extracted from post bodies) with the
Waller-Anderson community set \citep{WallerAnderson2021}. Restricting
edges to those with multiplicity at least $3$ and taking the largest
connected component gives a graph on $3{,}662$ subreddits and
$10{,}750$ edges, with average degree $5.87$. Every vertex carries
the Waller-Anderson partisan-axis score in $[-1, +1]$ as its feature.

Trained under EchoGAE's default hyperparameters on this graph, the
community-level Echo Chamber Score is $0.517$: within $0.02$ of the
$0.5$ no-discrimination midpoint, matching the pattern on Twitter.
The graph-blind ridge regression achieves $R^2 = 0.10$ on this
graph. The rewiring destroys the graph's clustering entirely
($0.291 \to 0.000$, all $26{,}413$ triangles removed). EchoGAE's
community score after retraining on the rewired graph is $0.504$: a
change of $-0.013$, comparable to the per-node residual noise. On
Reddit as on Twitter, EchoGAE's community score does not respond to
closure structure.

\subsection{Scope}

The empirical work now covers two estimators (EchoGAE and
InfoVGAE-lite) on five datasets (four Twitter, one Reddit). It does
not test TIMME \citep{XiaoTIMME2020}, which is a supervised
relational GCN. TIMME is a message-passing graph neural network and
therefore admits the same audit under Theorem~\ref{thm:main}. We defer that
work.

Two limitations of the Reddit audit deserve mention. First, the SNAP
soc-RedditHyperlinks dataset covers 2014--2017 hyperlinks between
subreddits, not the full user-comment interaction structure of
Reddit; the resulting subreddit graph is sparser and more
hyperlink-driven than a comment-based interaction graph would be.
Reconstructing a comment-based graph from Pushshift or ConvoKit
dumps is a heavier engineering effort we do not attempt. Second, the
Waller-Anderson partisan score is a projection of the same
$150$-dimensional community embeddings that many polarisation
estimators would themselves be trained to recover; using the score
as a label is a downstream-consistency check, not an independent
ground truth.
